% 付録再編草案・第2部：一歩の変化と微積分

\section{一歩の変化を細かくすると,微積分になるか}

直前では,更新の前後で「ずれ」がどう変わるかを調べた.
数列の変化を測る最も簡単な方法は,隣り合う項の差を取ることである.
\[
  \Delta a_n=a_{n+1}-a_n
\]
これを一階差分という.
たとえば $a_n=n^2$ なら,
\[
  \Delta a_n=(n+1)^2-n^2=2n+1
\]
となる.
二次式から一次式へ次数が一つ下がった.
微分が関数の次数を下げるのと同じように,差分は数列の変化を一段簡単にする.

差分を足すと,元の変化が戻る.
\[
  a_n-a_0
  =\sum_{k=0}^{n-1}(a_{k+1}-a_k)
  =\sum_{k=0}^{n-1}\Delta a_k
\]
である.
右辺では途中の項が打ち消し合い,最初から最後までの変化だけが残る.
微分の逆操作として積分があるように,差分の逆操作は和である.
この対応は,本文で階差型漸化式を解いた計算にもすでに現れていた.

今度は,一歩の大きさを明示しよう.
時間を幅 $h$ で区切り,\\,$t_n=nh$ における値を $a_n$ とする.
一単位時間あたりの変化は
\[
  \frac{a_{n+1}-a_n}{h}
\]
で測れる.
$h$ を小さくすると,差分商は微分係数に近づく.
逆に,微分方程式 $y'=F(t,y)$ の解を数列で近似するには,一歩の間の変化をおよそ $hF(t_n,a_n)$ と見積もればよい.
\[
  a_{n+1}=a_n+hF(t_n,a_n)
\]
とする方法がオイラー法である.
つまり,漸化式を細かい刻みで使えば微分方程式の計算になり,微分方程式を一歩ずつ計算すれば漸化式になる.

この対応を,最も簡単な微分方程式で確かめる.
\[
  \frac{dx}{dt}=\lambda x
\]
にオイラー法を使うと,
\[
  a_{n+1}=(1+h\lambda)a_n
\]
となり,等比数列が現れる.
一方,連続的な解は $x(t)=e^{\lambda t}x(0)$ である.
離散的な等比数列と連続的な指数関数は,同じ増加の規則を刻み幅の違う言葉で表している.
この例では,線形漸化式で特性根が成長や減衰を決めたのと同じように,微分方程式では $\lambda$ が成長や減衰を決める.

ただし,刻みを粗くしても連続的な動きが正しく再現されるとは限らない.
$\lambda<0$ なら微分方程式の解は減衰するが,オイラー法の数列が減衰するには
\[
  |1+h\lambda|<1
\]
が必要である.
刻み幅 $h$ が大きすぎると,元の連続的な解とは違って数列が振動したり発散したりする.
微分方程式を漸化式に直した後も,その漸化式の特性根を調べる必要がある.
離散と連続を結ぶと,両者の似ている点だけでなく,近似が崩れる条件も見えてくる.

微積分は,数列の近似計算だけでなく,漸化式そのものを設計するときにも役立つ.
方程式 $g(x)=0$ の根を求めたいとし,現在の近似値を $x_n$ とする.
$x_n$ の近くでは,関数を接線で近似できる.
小さな移動量 $h$ に対して
\[
  g(x_n+h)\approx g(x_n)+g'(x_n)h
\]
である.
接線上で関数の値を $0$ にするように $h$ を選ぶと,
\[
  h=-\frac{g(x_n)}{g'(x_n)}
\]
となる.
そこで次の近似値を
\[
  x_{n+1}=x_n-\frac{g(x_n)}{g'(x_n)}
\]
と定める.
これがニュートン法である.
方程式の根を一度に求めるのではなく,接線を使って根へ近づく更新規則を作った.

なぜこの更新は,根の近くで速く働くのだろうか.
根を $\alpha$ とし,誤差を $e_n=x_n-\alpha$ とする.
単純な根で関数が十分滑らかなら,テイラー展開によって
\[
  g(\alpha+e_n)
  =g'(\alpha)e_n+\frac12g''(\alpha)e_n^2+O(e_n^3)
\]
であり,導関数についても
\[
  g'(\alpha+e_n)=g'(\alpha)+g''(\alpha)e_n+O(e_n^2)
\]
となる.
この二つをニュートン法の式に代入すると,誤差の一次の項が打ち消され,
\[
  e_{n+1}
  =
  \frac{g''(\alpha)}{2g'(\alpha)}e_n^2
  +O(e_n^3)
\]
を得る.
普通の反復では誤差が定数倍に縮むのに対し,ニュートン法では十分近くで誤差がおよそ二乗になる.
一次近似によって更新を決め,テイラー展開によってその更新の精度を調べる.
微分と漸化式は,根へ近づく速さを理解するところで再び出会う.
